% !TeX root = ../../Book3.tex
\section{平面域的统一化}
\label{b3:sec:1204}

一般曲面的抛物型包括复环面，平面区域却不能容纳这个紧例外。还需回答一个看似较小的问题：一张与 \(\mathbb C^*\) 双全纯的区域，是否可能在平面中缺少许多点？不能仅凭它有一个孔就作判断；下面利用单叶函数在两端的奇点，把实际像确定下来。

\subsection{平面双曲区域}
\label{b3:subsec:120401}

\Needspace{6\baselineskip}
\begin{lemma}[去心平面的单叶平面像]\label{b3:lem:12-punctured-plane-embedding}
设 \(f:\mathbb C^*\to\mathbb C\) 为单射全纯函数。则存在 \(A\in\mathbb C^*\)、\(B\in\mathbb C\)，使
\[
 f(z)=Az+B\quad\hbox{或}\quad f(z)=\frac A z+B.
\]
特别地，\(f(\mathbb C^*)=\mathbb C\setminus\{B\}\)。
\end{lemma}
\begin{proof}
先排除0处的本性奇点。取 \(z_0\in\mathbb C^*\)，再取闭包避开0的一个小圆盘 \(V\ni z_0\)。开映射定理使 \(f(V)\) 包含 \(f(z_0)\) 的某个邻域。若0为本性奇点，Casorati--Weierstrass 定理给 \(z_n\to0\)，使 \(f(z_n)\to f(z_0)\)。充分大的 \(n\) 满足 \(z_n\notin V\)、\(f(z_n)\in f(V)\)，于是这个值既有 \(V\) 中的原像，又有 \(z_n\) 这个原像，违反单射。对 \(f(1/w)\) 在 \(w=0\) 重复相同论证，也排除了无穷远处的本性奇点。

因此两端都是可去奇点或极点。补入两端后得到球面亚纯函数，由\cref{b3:thm:03-sphere-rational}，它是有理函数 \(R=P/Q\)，其中 \(P,Q\) 互素。令 \(d=\max(\deg P,\deg Q)\)。可选一个有限值 \(c\)，避开 \(R(0)\)、\(R(\infty)\)、所有临界值，以及使 \(P-cQ\) 的最高次项消掉的至多一个值。方程 \(P-cQ=0\) 于是有 \(d\) 个互异的有限根；互素性保证这些根不是极点，而所选 \(c\) 保证它们均不为0。每个根都是原函数 \(f\) 的原像，单射性迫使 \(d=1\)。

所以 \(R\) 为 Möbius 变换。它在球面恰有一个极点，该点不能落在 \(\mathbb C^*\)，故只能是0或无穷远。极点为无穷远时 \(R=Az+B\)；极点为0时 \(R=A/z+B\)，且 \(A\ne0\)。两式的像直接给出结论。
\end{proof}
这里单射性既排除了本性奇点，也排除了高次有理映射。若只假设导数不为零，\(e^z\) 在无穷远仍有本性奇点；局部可逆不足以完成这一步。

\Needspace{6\baselineskip}
\begin{proposition}[平面区域的圆盘覆盖]\label{b3:prop:12-planar-hyperbolic}
采用深层统一化定理。设 \(U\subseteq\mathbb C\) 是非空连通开集。则 \(U\) 具有全纯圆盘万有覆盖，当且仅当 \(\mathbb C\setminus U\) 至少含两个不同点。其余情形恰为 \(U=\mathbb C\) 或 \(U=\mathbb C\setminus\{a\}\)，二者都具有平面万有覆盖。
\end{proposition}
\begin{proof}
先设补集至少有两个点。统一化给 \(U\) 的三种可能覆盖模型。球面模型的商只有球面，而非空平面区域不紧，故排除球面。若为平面模型，由前节的平面商分类，\(U\) 必与 \(\mathbb C\)、\(\mathbb C^*\) 或某个复环面双全纯。紧复环面仍被排除。

若存在双全纯映射 \(f:\mathbb C\to U\)，把它视为单叶整函数，\cref{b3:thm:05-entire-univalent}给 \(f(z)=Az+B\)，从而 \(U=\mathbb C\)。若存在双全纯映射 \(f:\mathbb C^*\to U\)，前一引理给 \(U=\mathbb C\setminus\{B\}\)。两种都与至少缺少两个点矛盾，因此只剩圆盘模型。

反向，补集为空时，恒等映射就是平面万有覆盖；补集只有 \(a\) 时，映射 \(w\mapsto a+e^w\) 是全纯覆盖，源平面单连通，故为平面万有覆盖。模型唯一性排除这两种区域再有圆盘万有覆盖。
\end{proof}
这是平面补集的计数；换到球面，还要把无穷远算作一个遗漏点。因此球面域的相应门槛是至少遗漏三个点。证明不要求补集有限，也不要求边界光滑或局部连通。只要选出两个不同的遗漏点，剩下的补集可以具有任意复杂形状。

\ProseParagraphBreak
内部圆盘覆盖与边界对应仍是不同结论。这里的投影在每个内部点附近都有完整的局部层，却不保证延伸为闭圆盘到普通闭包的连续映射。第七章关于边界与素端的条件仍须另行核对。

\subsection{特殊抛物型情形}
\label{b3:subsec:120402}

“真平面域”足以在单连通假设下使用 Riemann 映照，却不足以在任意区域上保证圆盘覆盖。\(\mathbb C^*\) 就是真平面域，它的指数覆盖是整个平面。与之对照，去心盘 \(\mathbb D^*\) 有上半平面覆盖，而有限圆环有条带覆盖。三者都是一个圆圈的同伦类型，复结构却区分了它们的端部。

\ProseParagraphBreak
删点的临界现象也可在球面上直接列出。不给球面删点，模型仍为球面；删去一点，用把该点送到无穷远的 Möbius 坐标得到平面；删去两点，可把它们送到 \(0,\infty\)，得到去心平面。删去至少三个点时，先把一个遗漏点送到无穷远，余下平面补集至少有两个点，前一命题遂给圆盘覆盖。

\ProseParagraphBreak
特别地，三穿孔球面可归一化为 \(\mathbb C\setminus\{0,1\}\)。统一化保证其圆盘覆盖存在，本章并未求出投影的显式公式。许多应用只需要源为圆盘、投影满射且局部双全纯；在这些应用中，缺少公式并不妨碍使用覆盖提升定理。下文的函数刚性就是一个例子。

\subsection{Poincaré 度量构造}
\label{b3:subsec:120403}

圆盘模型还带来一个自然的局部长度单位。由 Schwarz--Pick 的自同构等号，对任意圆盘自同构 \(T\)，
\[
 \frac{2|T'(w)|}{1-|T(w)|^2}=\frac2{1-|w|^2}.
\]
因而密度 \(\lambda_{\mathbb D}(w)\coloneqq2/(1-|w|^2)\) 乘上微小线段长度后，在圆盘坐标变换下不变。要把它传到底曲面，必须先检查两个选择：同一点的不同提升，以及底曲面的不同局部坐标。

\Needspace{6\baselineskip}
\begin{proposition}[圆盘长度元的下降]\label{b3:prop:12-poincare-descent}
设 \(X\) 为连通 Riemann 曲面，\(p:\mathbb D\to X\) 为全纯万有覆盖。在 \(X\) 的局部复坐标 \(z:V\to z(V)\subseteq\mathbb C\) 中，对 \(x\in V\) 任取 \(w\in p^{-1}(x)\)，定义
\[
 \lambda_z(z(x))=
 \frac{2}{(1-|w|^2)\,\bigl|(z\circ p)'(w)\bigr|}.
\]
这里导数在 \(w\) 附近映入 \(V\) 的覆盖小片上计算。这个正光滑密度与提升点、万有覆盖的选择无关；若另一坐标为 \(\zeta=\psi(z)\)，则
\[
 \lambda_\zeta(\psi(z))\,|\psi'(z)|=\lambda_z(z).
\]
因此各坐标中的 \(\lambda_z(z)|\mathrm dz|\) 拼成 \(X\) 上唯一使 \(p\) 局部保持圆盘长度元的共形长度元。
\end{proposition}
\begin{proof}
覆盖局部双全纯，故分母中的导数不为零。若另一个提升为 \(w_2\)，万有覆盖的甲板群在纤维上传递，所以存在圆盘自同构 \(T\)，满足 \(w_2=T(w)\)、\(p\circ T=p\)。在局部坐标中微分得
\[
 (z\circ p)'(w)=(z\circ p)'(T(w))\,T'(w).
\]
结合圆盘密度不变性，得到
\[
 \frac{\lambda_{\mathbb D}(T(w))}
      {|(z\circ p)'(T(w))|}
 =\frac{\lambda_{\mathbb D}(w)}
       {|(z\circ p)'(w)|}.
\]
于是密度与提升无关。固定一个局部逆 \(s:V\to\mathbb D\)，令 \(\sigma=s\circ z^{-1}\)，以 \(u\in z(V)\) 表示坐标数值，则同一公式写成
\[
 \lambda_z(u)=\frac{2|\sigma'(u)|}{1-|\sigma(u)|^2}.
\]
局部逆全纯且导数不为零，故所得密度光滑且正。

换成坐标 \(\zeta=\psi(z)\) 时，链式法则在分母多出 \(|\psi'(z)|\)，恰得陈述中的变换式，因而长度元相同。若改用另一个圆盘万有覆盖 \(p_1\)，覆盖唯一性给圆盘自同构 \(A\)，使 \(p_1\circ A=p\)。以刚才相同的链式法则及密度不变性比较两式，结果仍相同。最后，要求投影局部保持圆盘长度元，就必有
\[
 \lambda_z(z(p(w)))\,|(z\circ p)'(w)|=\lambda_{\mathbb D}(w),
\]
这逐点确定了密度，所以也给唯一性。
\end{proof}
该长度元称为 \term[poincare-duliang]{Poincaré 度量}[Poincaré metric]。它先定义分段光滑曲线的长度
\[
 L(\gamma)=\int \lambda_z(z(\gamma(t)))\,
                 |(z\circ\gamma)'(t)|\,\mathrm dt;
\]
曲线跨过坐标图时分段积分，坐标变换式保证结果一致。路径距离还需对连接两点的曲线长度取下确界。距离公式、完备性和曲率在第十三章独立展开，这里只完成下降本身。

\ProseParagraphBreak
抽象公式可以在两个覆盖上直接算出。先用 Cayley 变换 \(C(w)=(w-\mathrm i)/(w+\mathrm i)\) 把上半平面映到圆盘，有
\[
 1-|C(w)|^2=\frac{4\operatorname{Im}w}{|w+\mathrm i|^2},
 \quad |C'(w)|=\frac2{|w+\mathrm i|^2}.
\]
拉回圆盘长度元，便得上半平面密度 \(1/\operatorname{Im}w\)。这个系数与圆盘中的2对应，后面不再另作缩放。

\Needspace{6\baselineskip}
\begin{proposition}[去心盘与圆环的密度]\label{b3:prop:12-puncture-annulus-density}
设 \(\mathbb D^*=\{z:0<|z|<1\}\)，\(0<r<1\)，\(A(r,1)=\{z:r<|z|<1\}\)，并令 \(L=\log(1/r)\)。以圆盘密度 \(2/(1-|w|^2)\) 归一化的 Poincaré 密度分别为
\[
 \lambda_{\mathbb D^*}(z)=\frac1{|z|\log(1/|z|)},\quad
 \lambda_{A(r,1)}(z)=
 \frac{\pi}{L|z|\sin\!\left(\pi\log(1/|z|)/L\right)}.
\]
\end{proposition}
\begin{proof}
覆盖 \(p(w)=e^{\mathrm i w}\) 将上半平面映到去心盘，且
\[
 |p'(w)|=|p(w)|,\quad
 \operatorname{Im}w=\log(1/|p(w)|).
\]
将上半平面密度除以 \(|p'(w)|\)，得到第一式。

对圆环，使用条带 \(S=\{w:0<\operatorname{Im}w<L\}\) 及同一个公式 \(p(w)=e^{\mathrm i w}\)。映射
\[
 h(w)=e^{\pi w/L}
\]
把 \(S\) 双全纯映到上半平面：辐角在 \((0,\pi)\) 中，模与辐角分别确定实部和虚部。记 \(y=\operatorname{Im}w\)，则
\[
 \frac{|h'(w)|}{\operatorname{Im}h(w)}
 =\frac{\pi}{L\sin(\pi y/L)}.
\]
这就是条带密度。再除以 \(|p'(w)|=|z|\)，并以 \(y=\log(1/|z|)\) 代入，得到第二式。两个投影的纤维均为水平 \(2\pi\mathbb Z\) 平移轨道，条带和半平面都单连通，所以确实使用了全纯万有覆盖。
\end{proof}
密度只依赖 \(|z|\)，但并非常数；靠近不同端部的增长率也不同。去心盘径向曲线从半径 \(\rho_0\) 走到 \(\varepsilon<\rho_0\) 的长度为
\[
 \int_\varepsilon^{\rho_0}\frac{\mathrm d\rho}{\rho\log(1/\rho)}
 =\log\frac{\log(1/\varepsilon)}{\log(1/\rho_0)},
\]
随 \(\varepsilon\downarrow0\) 趋于无穷。这项计算展示穿孔在新长度中的位置；一般路径的距离与完备性还须由后章的完整论证处理。

\subsection{后续章节所需的统一化结论}
\label{b3:subsec:120404}

统一化的一个用途是把全纯映射提升到熟悉的函数空间。目标即使没有有界平面坐标，覆盖坐标也可能把映射化为有界整函数。

\Needspace{6\baselineskip}
\begin{proposition}[平面覆盖源到圆盘覆盖目标]\label{b3:prop:12-parabolic-to-hyperbolic}
设 \(X,Y\) 为连通 Riemann 曲面，\(q:\mathbb C\to X\)、\(p:\mathbb D\to Y\) 分别为全纯万有覆盖。则每个全纯映射 \(f:X\to Y\) 都恒常。
\end{proposition}
\begin{proof}
复合 \(f\circ q:\mathbb C\to Y\) 的源单连通。指定一个源点及其目标提升后，映射提升判别给
\[
 F:\mathbb C\to\mathbb D,\quad p\circ F=f\circ q.
\]
在覆盖小片中 \(F\) 是 \(p\) 的全纯局部逆与 \(f\circ q\) 的复合，因此为有界整函数。Liouville 定理使 \(F\) 恒常，故 \(f\circ q\) 恒常；\(q\) 满射，遂有 \(f\) 恒常。
\end{proof}
例如，一个整函数若遗漏两个不同有限值 \(a,b\)，就给从平面到 \(\mathbb C\setminus\{a,b\}\) 的全纯映射；后者由本节第一项命题有圆盘覆盖，故该整函数恒常。第十四章会把这一路线作为小 Picard 定理重新陈述，并与曲率方法比较。只遗漏一个值时，\(e^z\) 已是非常值反例，恰好落在平面覆盖目标 \(\mathbb C^*\) 上。

\ProseParagraphBreak
后续主线使用的输入可以准确分开。第十三章需要圆盘型曲面的全纯万有覆盖及甲板自同构不变性，然后由圆盘计算和路径提升建立距离理论；第十四章的上述省略值论证只需三穿孔球面的圆盘覆盖存在，不需要其显式模函数。选读章可以略过，但这些章会各自重述所需输入及来源。接受一般统一化、证明其具体后果以及计算一个模型，是三个可以分别检查的步骤。
